\section{The $K_4$ Vacuum Hypergraph}
\label{sec:hypergraph_model}

\subsection{Graph Construction}

We define a vacuum hypergraph $\mathcal{H} = (V, E)$ in the spirit of the Wolfram Physics Project.  The seed graph consists of:
\begin{itemize}
    \item A \textbf{$K_4$ complete graph} (``Oligon'') on 4 vertices $\{v_1, v_2, v_3, v_4\}$, with all $\binom{4}{2} = 6$ edges present.  This represents the minimal non-trivial complete graph with a non-degenerate eigenvalue spectrum.
    \item An \textbf{11-node vacuum ring} connecting nodes $v_5, \ldots, v_{15}$ in a cycle, with $v_5$ and $v_{15}$ attached to $v_1$ and $v_4$ respectively.
\end{itemize}
The resulting graph $\mathcal{H}$ has $|V| = 15$ nodes and $|E| = 17$ edges.  Its adjacency matrix $M \in \{0,1\}^{15\times 15}$ is symmetric with zero diagonal.

\subsection{Spectral Decomposition}
\label{subsec:spectral}

The spectral properties of $\mathcal{H}$ are dominated by the $K_4$ subgraph.  The adjacency matrix of $K_4$ alone has eigenvalues
\begin{equation}
    \mathrm{spec}(M_{K_4}) = \{3, -1, -1, -1\}\,,
    \label{eq:K4_eigenvalues}
\end{equation}
with spectral radius $\lambda_1 = 3$ (simple) and a triply degenerate eigenvalue $\lambda_2 = -1$.

The trace formula for powers of the $K_4$ adjacency matrix follows immediately:
\begin{equation}
    \mathrm{Tr}(M_{K_4}^n) = \sum_{i} \lambda_i^n = 3^n + 3(-1)^n\,,
    \label{eq:trace_formula}
\end{equation}
which corresponds to OEIS sequence A054878 (number of closed walks of length $n$ on $K_4$).  This has been verified computationally for $n = 1, \ldots, 10$ by explicit matrix exponentiation.

\subsection{The Spectral--Picard Bridge}
\label{subsec:spectral_picard}

The spectral radius $\lambda_1 = 3$ connects the discrete graph theory to algebraic geometry through the \emph{spectral--Picard bridge}.  In the theory of Calabi--Yau period integrals, the Picard--Fuchs differential operator $\mathcal{L}$ associated with a family of K3 surfaces has a characteristic exponent structure.  For the Cooper $s_{10}$ family~\cite{Cooper2012}, the recurrence
\begin{equation}
    u_n = \sum_{k=0}^{n} \binom{n}{k}^2 \binom{n+k}{k} \binom{2k}{k} \cdot (-4)^{n-k}
    \label{eq:cooper_recurrence}
\end{equation}
generates a sequence whose asymptotic growth rate is $|u_n| \sim C \cdot 3^n$, matching the spectral radius of $K_4$.  This coincidence is not accidental: the order-3 Picard--Fuchs operator $\mathcal{L}_3$ for Cooper $s_{10}$ has a symmetric square structure $\mathcal{L}_3 = \mathrm{Sym}^2(\mathcal{L}_2)$~\cite{AlmkvistVanStraten}, and the spectral radius of $\mathcal{L}_2$ equals $\sqrt{3}$, squaring to $\lambda_1 = 3$ under the symmetric product.

The Cooper $s_{10}$ K3 surface has Picard number $\rho = 19$ (formally verified in Lean~4 as theorem \texttt{spectral\_picard\_bridge}: $\lambda_1 = 3 \implies \rho = 19$).  Thus the $K_4$ graph \emph{deterministically selects} the Cooper $s_{10}$ K3 surface from the space of all K3 compactifications.

\subsection{Physical Interpretation: Oligon Topological Defects}

We interpret the $K_4$ clique as a \emph{topological defect} (``Oligon'') in the vacuum pregeometry.  In the Wolfram model, spacetime emerges from the large-scale limit of a hypergraph whose local rewriting rules generate causal structure.  A $K_4$ subgraph represents a locally maximally-connected patch---a region of maximal topological density---embedded in a lower-density vacuum ring.  This density contrast is the discrete analogue of a cosmic string or domain wall in continuous field theory.

The Oligon oscillates as the hypergraph evolves under its rewriting rule (Section~\ref{sec:hadamard_mask}), and these oscillations source gravitational waves via the quadrupole formula (Section~\ref{sec:gw_predictions}).  The frequency and amplitude of the gravitational radiation are entirely determined by the eigenvalue structure of $K_4$ and the physical scales assigned by the continuum limit (Section~\ref{sec:continuum_limit}).
